\chapter{Introduction}

\section{Problem Definition}
Banking institutions sit on large volumes of operational data: account histories, transactions, demographics, and loan records. The information that predicts whether a borrower will repay a loan is usually already present in that history before the loan is granted, but it is buried in millions of individual transactions and is not directly usable by a credit officer.

This project addresses the resulting decision problem as a supervised binary classification task. Given a loan application and the applicant's pre-loan transaction history, the system estimates whether the loan will be repaid (a \emph{good} loan) or will default (a \emph{bad} loan). The task is performed on the PKDD'99 Financial Dataset, which contains 682 loans and roughly one million account transactions from a Czech bank between 1993 and 1998. Of the 682 loans, 606 were repaid and 76 defaulted, giving a base default rate of 11.14 percent.

The central constraint is temporal: every behavioral feature is computed only from transactions dated strictly before the loan grant date. Post-loan activity is never used, so the model reproduces what the bank could have known at decision time rather than what it learned afterwards.

\section{Motivation and Business Context}
Loan defaults are direct losses, while rejected good customers are opportunity costs. A model that separates the two groups well allows the bank to route applications into a review queue, ask for additional collateral, or adjust the offered amount. The cost of the two error types is not symmetric: approving a future defaulter is usually more expensive than reviewing a good customer twice.

The dataset is also strongly imbalanced. With only 11.14 percent positives, a trivial classifier that predicts ``repaid'' for every application reaches 88.86 percent accuracy while catching zero defaulters. Accuracy is therefore not a useful headline metric here, and the evaluation has to focus on the positive class: precision, recall, F1-score, Average Precision, and ROC-AUC. This project treats that methodological point as part of the problem, not as an afterthought.

\section{Project Objectives}
The project covers the complete data mining pipeline required by Units III and IV of the syllabus. Its objectives are:

\begin{enumerate}
    \item To clean and integrate the raw PKDD'99 tables into a single loan-level analytical dataset.
    \item To engineer behavioral features from pre-loan transactions only, and to verify that no post-loan information leaks into the feature set.
    \item To perform exploratory data analysis and visualize the behavioral differences between repaid and defaulted loans.
    \item To train and compare six classification algorithms: Decision Tree, Naive Bayes, Random Forest, Support Vector Machine (SVM), K-Nearest Neighbors (KNN), and Gradient Boosting.
    \item To handle class imbalance explicitly by training each classifier twice, once on the raw data and once on a SMOTE-balanced training set.
    \item To tune hyperparameters and select the final model using repeated stratified cross-validation rather than a single train--test split.
    \item To interpret the selected model through feature importance and association rule mining, and to translate the findings into business recommendations.
\end{enumerate}

\section{Scope and Deliverables}
The project scope is limited to binary loan-default classification on historical data. It does not include live scoring, credit-policy design, or causal analysis. The supplementary components are customer segmentation with KMeans, association rule mining with Apriori and FP-Growth, and a Streamlit decision-support console that reads the exported result files.

Table~\ref{tab:deliverables} maps each expected deliverable to the artifact that implements it.

\begin{table}[H]
    \centering
    \small
    \renewcommand{\arraystretch}{1.25}
    \begin{tabularx}{\textwidth}{|>{\raggedright\arraybackslash}p{4.6cm}|>{\raggedright\arraybackslash}X|}
        \hline
        \textbf{Deliverable} & \textbf{Where it is implemented} \\
        \hline
        Problem definition and objectives & This chapter and \texttt{docs/Problem\_Statement.md} \\
        \hline
        Dataset description and source & Chapter 2; raw CSVs in \texttt{data/} \\
        \hline
        Data preprocessing pipeline & Chapter 4; \texttt{01\_preprocessing.ipynb}, \texttt{02\_features.ipynb} \\
        \hline
        EDA and visualizations & Chapter 5; \texttt{03\_eda.ipynb}; 8 EDA figures, 14 project figures in total \\
        \hline
        Classification algorithms & Chapter 6; \texttt{04\_modeling.ipynb} \\
        \hline
        Performance comparison & Chapter 6; \texttt{metrics\_comparison.csv}, \texttt{tuning\_results.csv}, \texttt{cv\_comparison.csv} \\
        \hline
        Conclusions and recommendations & Chapter 7; \texttt{docs/Jury2\_Report.md} \\
        \hline
        Association rule mining (Unit III) & Section 6.8; \texttt{tools/pattern\_mining.py} \\
        \hline
        Project report and presentation & This document and the executed notebooks \\
        \hline
    \end{tabularx}
    \caption{Deliverable coverage for Jury II}
    \label{tab:deliverables}
\end{table}

\section{Jury II Evaluation Coverage}
\label{sec:evaluation-coverage}

Table~\ref{tab:jury2-coverage} maps the evaluation parameters for Jury II to the parts of this report that address them. The table is an index, not a substitute for the chapters: each row points to the analysis, table, or figure that contains the evidence.

\begin{table}[htbp]
    \centering
    \small
    \renewcommand{\arraystretch}{1.2}
    \begin{tabularx}{\textwidth}{@{}>{\raggedright\arraybackslash}p{4.4cm}X@{}}
        \toprule
        \textbf{Evaluation parameter} & \textbf{Coverage in this report} \\
        \midrule
        Dataset & Chapter~\ref{ch:literature}, Table~\ref{tab:pkdd_tables}; source joins and cleaning in Chapter~\ref{ch:preprocessing}. \\
        \addlinespace
        Architecture diagram & Figure~\ref{fig:architecture} and the stage description in Section~\ref{sec:pipeline-stages}. \\
        \addlinespace
        Input / output & Section~\ref{sec:io}, Table~\ref{tab:io}; the 18-feature model input and scored output are stated there. \\
        \addlinespace
        Classes / labels & Section~\ref{sec:target-definition}, Table~\ref{tab:target-balance}; the A/C versus B/D mapping and the 606/76 balance. \\
        \addlinespace
        Pattern matching (FP-Growth, Apriori) & Section~\ref{sec:pattern-mining}, Table~\ref{tab:pattern-rules}; both algorithms, the 41 itemsets, and the two default rules. \\
        \addlinespace
        Comparison with other classification algorithms & Sections~\ref{sec:test-results} and~\ref{sec:cv-results}, Tables~\ref{tab:test-results} and~\ref{tab:cv-results}; six classifiers under baseline and SMOTE. \\
        \addlinespace
        Classification algorithm & Chapter~\ref{ch:literature} and Sections~\ref{sec:classifiers} and~\ref{sec:training-pipeline}; the six learners and their pipelines. \\
        \addlinespace
        Innovation added to the classification algorithm & Section~\ref{sec:innovation}; the pre-loan representation, balance-capacity view, two-arm validation, and rule-plus-model review layer. \\
        \addlinespace
        Performance metrics & Section~\ref{sec:evaluation-protocol} and Sections~\ref{sec:test-results} and~\ref{sec:ranking-curves}; accuracy, precision, recall, F1, AP, ROC-AUC, and confusion matrices. \\
        \addlinespace
        Chart & Chapter~\ref{ch:eda} and Figures~\ref{fig:roc-curves} and~\ref{fig:pr-curves}; the Streamlit console in Section~\ref{sec:streamlit} presents the same charts. \\
        \bottomrule
    \end{tabularx}
    \caption{Coverage of the Jury II evaluation parameters.}
    \label{tab:jury2-coverage}
\end{table}

\section{Organization of the Report}
Chapter 2 reviews the dataset, the classifiers, imbalance handling, and the evaluation metrics. Chapter 3 describes the system architecture, the input and output specification, leakage controls, and the evaluation protocol. Chapter 4 documents data cleaning and feature engineering. Chapter 5 presents the exploratory analysis. Chapter 6 reports the model implementation, tuning, comparison, feature importance, and pattern mining. Chapter 7 states the conclusions, innovation, limitations, and future scope. Section~\ref{sec:evaluation-coverage} maps each Jury II evaluation parameter to these chapters.
